\section{The LASM Framework}
\label{sec:framework}

The Large-scale Adversarial Security Metric (LASM) framework is designed to bridge the severe theoretical gap between static, isolated model evaluation and dynamic, system-level enterprise deployment. Unlike existing paradigms that treat Generative AI components as standalone text prediction engines, LASM models the entire socio-technical system iteratively. To accomplish this, we formalize five continuously interacting components that collectively map technical exploitability, temporal decay, and game-theoretic deterrence.

\subsection{LLM Attack Surface (LAS)}
\label{subsec:las}

We abandon naive combinatorial attack surface metrics in favor of stochastic reachability analysis via Absorbing Markov Chains (AMC). An LLM-based system is modeled as a state space $\mathcal{X} = \mathcal{V} \cup \mathcal{A}$, where $\mathcal{V}$ are transient states representing discrete processing nodes (e.g., the LLM orchestrator, a RAG retrieval engine, external API executors) and $\mathcal{A}$ comprises absorbing states such as $s_{comp}$ (successful malicious compromise) and $s_{safe}$ (safety termination).

The transition matrix $P$ captures information flow carrying malicious permutations between nodes. We block $P$ into $Q$ (transient probabilities) and $R$ (absorption probabilities). This allows us to compute the \textbf{Fundamental Matrix} $N = (I - Q)^{-1}$, mapping all feedback loops and multi-hop interactions intrinsically.

The exact expected exploitability, defining our LLM Attack Surface ($LAS$) under a computational budget $B$ at deployment time $t$ for target system $S$, is defined as:
\begin{equation}
    LAS(S, t, B) = h^T \cdot (I - Q_S)^{-1} \cdot P_{comp}(t, B)
\end{equation}

Here, $P_{comp}(t, B)$ is the empirical exploitation emission vector measuring independent component failures, and $h$ assigns a harm severity coefficient based on standardized taxonomies (e.g., MITRE ATLAS). 

By adopting the AMC formulation, we intrinsically calculate complex multi-agent cascading failures without relying on arbitrary, additive composability scalars, thereby solving the non-linear interaction gap mathematically.

\subsection{Standardized Attack Success Rate (SASR)}
\label{subsec:sasr}

Once the attack surface is mapped, we must accurately measure the discrete vulnerability. The prevailing metric, binary Attack Success Rate (ASR), is fundamentally flawed; it frequently succumbs to the \textit{False-Safety Paradigm}, granting passing scores to "illusionary safe" responses such as compliant gibberish or partial disclosures. 

To rectify this reporting bias, LASM introduces the Standardized Attack Success Rate (SASR), a continuous vulnerability metric computed over a standard adversarial prompt set $Q$ subject to a Normalized Query Unit (NQU) Budget $B$. To guarantee metric credibility, SASR employs a rigorous \textbf{Three-Judge Protocol (3JP)}:

\begin{enumerate}
    \item \textbf{Heuristic Classifier ($C_1$):} A rapid, rule-based classifier that filters out absolute, hard-coded safety refusals.
    \item \textbf{High-Capacity LLM Judge ($C_{LLM}$):} A frontier model (e.g., \texttt{Qwen2.5-3B}) instructed with specialized rubrics to continuously score the \textit{specificity}, \textit{convincingness}, and \textit{actionability} of the adversarial execution.
    \item \textbf{Human-in-the-Loop Validation:} A mandatory threshold ensuring inter-rater reliability via Cohen's Kappa ($k \ge 0.6$).
\end{enumerate}

For any given adversarial query $q \in Q$, the localized SASR is computed as:
\begin{equation}
    SASR_{q} = \frac{SR_{C_1}(q) + SR_{C_{LLM}}(q)}{2}
\end{equation}
By aggregating $SASR_q$ across $Q$, defenders extract a substantially more realistic penetration probability that resists naive evasions.

\subsection{Temporal Robustness Curve (TRC)}
\label{subsec:trc}

Security in Generative AI is not a static property but a continuously decaying state. As novel jailbreaking techniques (e.g., GCG, PAIR) are open-sourced, they trigger a "cascade" of derivative attacks. To model this, the Temporal Robustness Curve (TRC) forecasts degradation using self-exciting \textbf{Hawkes Processes}.

The TRC evaluates the deployed defense mechanism $D$ against the mathematically strongest known adversarial attack algorithm $A^*$ at progressing time intervals $t$:
\begin{equation}
    TRC_{prospective}(D, t_i) = \mathbb{E}_{Hawkes}[SASR(A_i, D, Q, B_{mid})]
\end{equation}

By simulating the self-exciting temporal intensity $\lambda(t) = \mu + \sum_{t_i < t} \alpha e^{-\beta(t-t_i)}$, we mathematically forecast the \textbf{Defense Half-Life (DHL)}. The DHL indicates when a mechanism's penetration probability crosses the critical threshold, shifting enterprise security from reactive patching to lifecycle forecasting.

\subsection{Game-Theoretic Extension: Economic Attack Surface (EAS)}
\label{subsec:eas}

High technical vulnerability ($LAS$) rarely translates to guaranteed exploitation if the cost of attack is prohibitive. LASM formalizes deterrence using Behavioral Game Theory, abandoning rigid deterministic limits in favor of a Quantal Response Equilibrium (QRE).

We model boundedly rational adversaries across profiles $\pi$ (e.g., script-kiddies to nation-state APTs) using a continuous Logit choice function. The Economic Attack Surface ($EAS$) replaces the boolean indicator entirely:
\begin{equation}
    EAS(S, t, B, \pi) = h^T \cdot (I - Q_S)^{-1} \cdot \text{diag}(P_{attack}(\pi)) \cdot P_{comp}(t, B)
\end{equation}
where the probability of the adversary executing the attack is intrinsically weighted by their rationality parameter $\lambda_\pi$:
\begin{equation}
    P_{attack}(\pi)_i = \frac{1}{1 + \exp\Big(-\lambda_\pi \cdot \big( \mathbb{E}[G_{A^*}] - C_{A^*} - \varepsilon_\pi \big) \Big)}
\end{equation}

This differentiable formulation ensures that incremental investments in friction continuously limit probabilistic exploit behavior, allowing defenders to utilize convex optimization for strict security budgeting.

\subsection{Attack Deterrence Threshold (ADT)}
\label{subsec:adt}

Building directly upon the EAS, defenders historically struggle with the \textit{over-defense problem}—stacking expensive, high-latency safety filters long after the threat model has been deterred. LASM provides a formal stopping criterion via the Attack Deterrence Threshold (ADT).

The ADT identifies the exact breakpoint for the optimal defense investment multiplier $\delta^*$:
\begin{equation}
    \delta^* = \frac{V_{target} - C_{opp}}{C_{compute}}
\end{equation}
When an enterprise scales their defense mechanism such that $\delta > \delta^*$, the Expected ROI for the adversary becomes mathematically negative. Under these conditions, the system secures an effective 0\% EAS with $100\%$ economic deterrence. The ADT ensures that security teams allocate resources rationally relative to their specific threat environment, establishing LASM as the foundation for modern enterprise AI security management.
